\documentclass[12pt]{article}
\usepackage[utf8]{inputenc}
\usepackage[T1]{fontenc}
\usepackage{hyperref}
\title{Gbe, Weme Language Resources for K-12 Curriculum Development}
\author{Compiled by Claude Sonnet 5 (claude-sonnet-5) for Andr\'as Kornai}
\date{}
\begin{document}
\maketitle
\section{Introduction}
Weme Gbe (ISO 639-3 wem; also called Ouem\'e) is spoken in Benin. It is a Gbe language, grouped with Fon, Gun and related varieties. The search found no resources specific to Weme. The results cover only the neighbouring Fon and Gun languages.
\section{Primary Educational Resources}
None found for Weme.
\section{Linguistic and Reference Materials}
The following concern related languages, not Weme itself, and are listed only as context:
\begin{itemize}
\item Gun language overview: \url{https://en.wikipedia.org/wiki/Gun_language}
\item Fon language overview: \url{https://en.wikipedia.org/wiki/Fon_language}
\item Standard uniform orthography of Gbe (Capo's proposal): \url{https://fr.wikipedia.com/wiki/Orthographe_uniforme_standard_du_gbe}
\item Fon--French dictionary (Fon, not Weme): \url{https://www.koeppe.de/titel_dictionnaire-fon-francais-francais-fon}
\end{itemize}
\section{Religious and Community Resources}
None found.
\section{Digital and Audio Resources}
None found.
\section{Curriculum Development Context}
No Weme-specific orthography, dictionary or school material turned up. A curriculum might build on the unified Gbe orthography and Fon materials, but whether these suit Weme would need to be checked with Weme speakers and Beninese language authorities.
\end{document}